\documentclass{article}
\usepackage{amsmath,amssymb,amsthm,mathtools}
\newtheorem{theorem}{Theorem}
\newtheorem{lemma}{Lemma}
\newtheorem{proposition}{Proposition}
\newtheorem{corollary}{Corollary}
\newtheorem{definition}{Definition}
\newtheorem{example}{Example}
\title{ShadowBench analysis/L2/ana\_four\_L2\_004}
\begin{document}
\maketitle
% Problem id: analysis/L2/ana_four_L2_004
% Companion instructions: docs/instructions.md
% Candidate skeletons: docs/skeletons/
% Lean target: ShadowBench/Source/Main.lean

\begin{theorem}[MeasureTheory.Integrable.fourierInv_fourier_eq] Let $f$ be an integrable function on a finite-dimensional real inner product space. If its Fourier transform $\mathcal{F} f$ is also integrable, then at every point where $f$ is continuous, the inverse Fourier transform of $\mathcal{F} f$ equals $f$ itself; that is, $\mathcal{F}^{-1}(\mathcal{F} f)(v) = f(v)$ for all continuity
 points $v$ of $f$.
\end{theorem}
\end{document}
